\chapter{Fasci}

Sia \(X\) uno spazio topologico. Allora l'insieme \(\mathrm{Open}{(X)}\) degli
aperti di \(X\), ordinato da \(\supseteq\), forma una categoria
\(\mathrm{Open}{(X)}\).

\section{Prefasci e fasci}

\begin{definition}{Prefasci}
    La categoria dei \emph{prefasci} su \(X\) a valori in una categoria
    \(\mathcal{C}\) è 
    \[
      \mathtt{Pfh}{(X, C)} = \mathtt{Fun}{(\mathrm{Open}{(X)}^{op}, \mathcal{C})}
    \]
    cioè \(\mathcal{F} \in \mathtt{Pfh}{(X, \mathcal{C})}\) è dato da
    \(\mathcal{F}{(U)} \in \mathcal{C}\) per ogni \(U \subseteq X \) aperto e da
    un morfismo di \(\mathcal{C}\), \(\rho_{V,U} : \mathcal{F}{(V)} \to
    \mathcal{F}{(U)}\) per ogni \(U \subseteq V \subseteq X  \) aperti tale che
    \(\rho_{U,U} = 1_{\mathcal{F}{(U)}} \) e \(\forall U \subseteq V \subseteq W
    \) aperti, \(\rho_{W, U} = \rho_{V, U} \circ \rho_{W, V} \). Inoltre per
    ogni \(\mathcal{F}, \mathcal{G} \in \mathtt{Pfh}{(X, \mathcal{C})}\), un
    morfismo \(\varphi  : \mathcal{F} \to \mathcal{G}\) è dato da \(\varphi_{U}
    : \mathcal{F}{(U)} \to \mathcal{G}{(U)}\) morfismo di \(\mathcal{C}\) tale
    che \(\forall U \subseteq V \subseteq X  \) aperti, 
    \[
      \varphi_U \circ \rho_{V,U} ^{\mathcal{F}} = \rho_{V, U} ^{\mathcal{G}}
      \circ \varphi_V
    \]
    o equivalentemente il seguente diagramma commuta:
\[\begin{tikzcd}
	{\mathcal{F}{(V)}} & {\mathcal{F}{(U)}} \\
    {\mathcal{G}{(V)}} & {\mathcal{G}{(U)}}
	\arrow["{\rho_{V,U} ^{\mathcal{F}}}", from=1-1, to=1-2]
	\arrow["\varphi_V", from=1-1, to=2-1]
	\arrow["\varphi_U", from=1-2, to=2-2]
	\arrow["\rho_{V,U} ^{\mathcal{F}}"', from=2-1, to=2-2]
\end{tikzcd}\]
    e per ogni \(U \subseteq V \subseteq X  \) aperti, data \(\sigma \in
    \mathcal{F}{(V)}\), invece di \(\rho_{V,U} {(\sigma)} \in \mathcal{F}{(U)}\)
    si scrive \(\sigma_{U}\) 
\end{definition}
D'ora in poi assumiamo \(\mathcal{C}\) concreta (in particolare \(\mathtt{Set},
{(\mathtt{Grp})}, \mathtt{Rng}, A\mathtt{-Mod}\))

\begin{definition}{Fascio}
    \(\mathcal{F} \subseteq \mathtt{Pfh}{(X, C)} \) è un \emph{fascio} se
    \(\forall U \subseteq \mathrm{Open}{(X)} \) e per ogni \(\{U_{i}: i \in I\}
    \) ricoprimento aperto di \(U\) (cioè \(U_{i} \in \mathrm{Open}{(X)}\), per
    ogni \(i \in I\) e \(U = \bigcup_{i \in I} U_{i}\))).

    Dati \(\sigma_{i} \in \mathcal{F}{(U_{i})}\) (\(\forall i \in I\)) tali che
    \(\sigma_{i} |_{U_{i} \cap U_{j}} = \sigma_{j} |_{U_{i} \cap U_{j}} \) per
    ogni \(i, j \in I\), allora \(\exists ! \, \sigma \in \mathcal{F}{(U)}\)
    tale che \(\sigma_{i} = \sigma|_{U_{i}} \) per ogni \(i \in I\).

    Se Vale solo l'unicità, si dice che il prefascio \(\mathcal{F}\) è
    \emph{separato}.
\end{definition}

Siano \(\mathcal{F}, \mathcal{G}\) due fasci. Un morfismo \(\mathcal{F} \to
\mathcal{G}\) di fasci è un morfismo di prefasci. 

\(\mathtt{Sh}{(X, \mathcal{C})}\) è la sottocategoria piena di
\(\mathtt{PSh}{(X, \mathcal{C})}\) con oggetti i fasci.

\begin{note}[zione]
    Invece di \(\mathcal{F}{(U)}\) si può scrivere \(\Gamma_f{(U, % TODO forse no
    \mathcal{F})}\). Gli elementi di \(\mathcal{F}{(U)}\) si dicono anche
    \emph{sezioni} di \(\mathcal{F}\) su \(U\).
\end{note}

\begin{example}{}
    Sia \(S \in \mathcal{C}\). \(\mathcal{F}_{S} \in \mathtt{Sh}{(X,
    \mathcal{C})}\) definito da \(\mathcal{F}_S{(U)} := S^{U}\) con \(\rho_{V,U}
    : \mathcal{F}_S {(V)} \to \mathcal{F}_S{(U)}\) la vera restrizione di
    funzioni. Questo è un fascio perché, date \(\sigma_{i} : U_{i} \to S\) tali
    che \(\sigma_{i} |_{U_{i} \cap U_{j}} = \sigma_{j} |_{U_{i} \cap U_{j}}  \),
    allora esiste unica \(\sigma : U \to S\) tale che \(\sigma|_{U_{1}} =
    \sigma_{i}\) per ogni \(i \in I\), se \(\{U_{i}\}_{i \in I} \) è un
    ricoprimento aperto di \(U\).

    Analogamente, se \(S\) è anche uno spazio topologico (tali che le operazioni
    di gruppo / anello siano continue, ad esempio \(S = \mathbb{R}\) e \(S =
    \mathbb{C}\) con \(\mathcal{C} = \mathtt{Rng}\)). Posso allora considerare
    il fascio \(\mathcal{G}_{S} \) definito da \(\mathcal{G}_S {(U)} := \{G \in
    S^{U} : G\text{ è continua}\}\).

    Ci sono varianti di questo esempio se \(X\) è una varietà
    \emph{differenziabile}, \emph{complessa} o \emph{algebrica}. In tal caso
    si considerano funzioni \(C^{\infty}\), \emph{olomorfe} o
    \emph{regolari}.
\end{example}

\begin{remark}{}
    Sia \(\mathcal{F} \in \mathtt{PSh}{(X, \mathcal{C})}\). Se \(\mathcal{F}\) è
    \emph{separato}/\emph{fascio} allora \(\# \mathcal{F}{(\varnothing)} \le 1 /
    = 1\). Posso infatti prendere come ricoprimento aperto il ricoprimento
    vuoto, cioè costituito da nessun aperto, di \(\varnothing\), ottengo che
    \emph{è unico}/\emph{esiste unico} \(\sigma \in \mathcal{F}{(\varnothing)}\).
\end{remark}

\begin{example}{}
    Sia \(S \in \mathcal{C}\) tale che \(\# S > 1\). Allora il prefascio
    costante di valore \(S\) (cioè \(\mathcal{F}{(U)} = S\) per ogni \(U \in
    \mathrm{Open}{(X)}\) o \(\rho_{V, U} = 1_S\) per ogni \(U \subseteq V
    \subseteq X\)) non è separato.

    Il fascio (localmente) costante di valora \(S\) è \(S_X\) definito da \(S_X
    {(U)} = \{\sigma \in S^{U} : \sigma \text{ è localmente costante}\}\), dove
    localmente costante significa che \(\sigma\) è continua mettendo su \(S\) la
    topologia discreta (quindi la controimmagine di ogni sottoinsieme di \(S\) è
    un aperto)
\end{example}

\begin{definition}{Spazio anellato}
    Uno \emph{spazio anellato} è una coppia \({(X, O_X)}\) con \(X\) spazio
    topologico e \(O_X \in \mathtt{Sh}{(X, \mathtt{Rng})}\) (di solito si
    considerano anelli commutativi).
\end{definition}

\section{Fasci di moduli}

\begin{definition}{}
    Sia \({(X, O_X)}\) uno spazio anellato. Un prefascio di \(O_X\)-moduli è un
    (pre)fascio \(\mathcal{F}\) di gruppi abeliani con una struttura di
    \(0_X{(U)}\)-modulo su \(\mathcal{F}{(U)}\) per ogni \(U\) aperto tale che
    \(\forall  U \subseteq V \subseteq X  \) aperti, \(\forall \sigma \in
    \mathcal{F}{(V)}\) e per ogni \(\alpha \in O_X{(V)} \), \((\alpha
    \sigma)|_{U}  = \alpha_U \sigma|_U\).
\end{definition}

SIano \(\mathcal{F}, \mathcal{G}\) (pre)fasci di \(O_X\)-moduli. Un morfismo
\(\varphi  : \mathcal{F} \to \mathcal{G}\) è un morfismo in \(PSh{(X,
\mathtt{Ab})}\) tale che \(\varphi _U : \mathcal{F}{(U)} \to \mathcal{G}{(U)}\)
è \(O_X{(U)}\)-lineare per ogni \(U \in \mathrm{Open}{(X)}\). Si ottiene una
cateogria \(O_X\mathtt{-PMod}\) dei prefasci di \(O_X\)-moduli con
sottocategoria piena \(O_X\mathtt{-Mod}\) dei fasci di \(O_X\)-moduli.

\begin{remark}{}
    So che \(\mathtt{PSh}{(X, \mathtt{Ab})}\) è una categoria abeliana (perché
    lo è \(\mathtt{Ab}\). Allora anche \(O_X\mathtt{-PMod}\) è abeliana
    (esercizio). 
\end{remark}

\begin{example}{}
    Sia \(X \neq \varnothing\) con la topologia concreta. Allora \(O_X\) è fascio
    di anelli su \(X\), allora \(O_X{(\varnothing)} = \{0=1\} \) e \(A =
    O_X{(X)}\) è un anello qualunque.
\end{example}

Se \(\mathcal{F} \in O_X\mathtt{-Mod}\), allora \(\mathcal{F}{(\varnothing)} =
O\) e \(\mathcal{F}{(X)}\) è un \(A\)-modulo qualunque. Segue che
\(O_X\mathtt{-Mod} \cong A\mathtt{-Mod}\). 

\begin{example}{}
    Sia \(A\) un anello. Allora \(A_X\) è un fascio di anelli su \(X\) tale che
    \(A_X\mathtt{-Mod}\cong \mathtt{Sh}{(X, A\mathtt{-Mod})}\).

    Infatti sia \(\mathcal{F} \in  A_X\mathtt{-Mod}\). Allora \(\forall U \in
    \mathrm{Open}{(X)}\), \(\mathcal{F}{(U)}\) è \(A_X{(U)}\)-modulo. Infatti
    considerando l'omomorfismo di anelli \(A \to A_X{(U)}\) dato da \(a \mapsto
    {(x \mapsto a)}\), per restrizione degli scalari \(\mathcal{F}{(U)}\) è
    anche \(A\)-modulo. Viceversa, se \(\mathcal{F} \in \mathtt{Sh}{(X,
    A\mathtt{-Mod})}\) devo dimostrare che \(\forall U \in
    \mathrm{Open}{(X)}\), \(\mathcal{F}{(U)}\) è anche \(A_X{(U)}\)-modulo.

    Data \(\sigma \in \mathcal{F}{(U)}\) e \(\alpha \in  A_X{(U)}\) devo
    definire \(\alpha \sigma \in  \mathcal{F}{(U)}\). \(U = \coprod_{a \in A}
    \alpha ^{-1}{(a)}\) ricoprimento aperto di \(U\) (perché \(\alpha\) è
    localmente costante), allora \(\exists ! \, \alpha \sigma \in
    \mathcal{F}{(U)}\) tale che \({(\alpha \sigma)}|_{\alpha^{-1}{(a)}} = a
    \sigma|_{\alpha^{-1}{(a)}}  \) .
\end{example}

\section{Morfismi e fascio associato}

\begin{definition}{}
    Sia \(\varphi : \mathcal{F}\to \mathcal{G}\) un morfismo di prefasci (in
    \(\mathtt{PSh}{(X, \mathcal{C})}\) o in \(O_X\mathtt{-PMod}\)). Allora
    \(\varphi \) è \emph{iniettivo}/\emph{suriettivo} se \(\varphi _U\) è
    \emph{iniettivo}/\emph{suriettivo} per ogni \(U \in \mathrm{Open}{(X)}\)
    (\(\varphi \text{biunivoco} \iff \varphi \text{iso}\)).

    \(\varphi \) è localmente iniettiva se \(\forall U \in  \mathrm{Open}{(X)}\)
    e \(\forall  \sigma, \sigma' \in  \mathcal{F}{(U)}\) tale che \(\varphi
    {(\sigma)} = \varphi {(\sigma')}\), allora \(\exists  \{U_{i} : i \in I\} \)
    ricoprimento aperto di \(U\) tale che \(\sigma|_{U_{i}} = \sigma'|_{U_{i}}
    \) per ogni \(i \in I\).

    \(\varphi \) è \emph{localmente suriettivo} se \(\forall U \in
    \mathrm{Open}{(X)}\) e \(\forall \tau \in \mathcal{G}{(U)}\), esiste un
    ricoprimento aperto di \(U\) \(\{U_{i}\}_{i \in I}  \) tale che
    \(\tau|_{U_{i}} \in \mathrm{im}{(\varphi _{U_{i}} )}\) per ogni \(i \in I\) 
\end{definition}

\begin{remark}{}
    Se \(\varphi \) è \emph{iniettivo}/\emph{suriettivo}, allora \(\varphi \) è
    localmente \emph{iniettivo}/\emph{suriettivo}.

    Se \(\varphi \) è localmente iniettivo e \(\mathcal{F}\) è separato, allora
    \(\varphi \) è iniettivo.

    % nota non suriettivo
\end{remark}

\begin{proposition}{}
    Sia \(\varphi : \mathcal{F} \to \mathcal{F}'\) un morfismo localmente
    biunivoco di prefasci, e sia \(\psi : \mathcal{F} \to \mathcal{G}\) un
    morfismo di prefasci con \(\mathcal{G}\) fascio. Allora esiste unico \(\psi'
    : \mathcal{F}^{\prime} \to \mathcal{G}\) morfismo di prefasci tale che
    \(\psi = \psi' \circ \varphi \) 
\end{proposition}

\begin{corollary}{}
    Siano \(\varphi_i : \mathcal{F} \to \mathcal{F}_i\) morfismi localmente
    biunivoci (per \(i=1,2\)) con \(\mathcal{F}_i\) fasci. Allora esiste unico
    \(\varphi : \mathcal{F}_1 \to \mathcal{F}_2\) tale che \(\varphi_{2} =
    \varphi  \circ \varphi_{1}\) e \(\varphi \) è isomorfismo.
\end{corollary}

\begin{proposition}{}
    Per ogni prefascio \(\mathcal{F}\) esiste \(\eta = \eta_\mathcal{F} :
    \mathcal{F} \to \mathcal{F}^{\mathcal{A}}\) localmente biunivoco con
    \(\mathcal{F}^{\mathcal{A}}\) fascio (detto \emph{fascio associato a}
    \(\mathcal{F}\))
\end{proposition}

\begin{corollary}{}
    Esiste il \emph{funtore associato}
    \[
      {(-)}^{\mathcal{A}} : \mathtt{PSh}{(X, \mathcal{C})} \to \mathtt{Sh}{(X,
      \mathcal{C})}
    \]
    (oppure \({(-)}^{\mathcal{A}} : 0_X\mathtt{-PMod} \to 0_X\mathtt{-Mod}\))
    aggiunto sinistro dell'inclusione \(\mathtt{Sh}{(X, \mathcal{C})}
    \hookrightarrow \mathtt{Psh}{(X, \mathcal{C})}\) (o rispettivamente
    \(0_X\mathtt{-Mod} \hookrightarrow 0_X\mathtt{-PMod}\)).

    Inoltre per ogni \(\mathcal{F}\) prefascio, \(\eta_\mathcal{F} : \mathcal{F}
    \to \mathcal{F}^{\mathcal{A}}\) è la composizione dell'unità
    dell'aggiunzione
\end{corollary}

\begin{proof}[Dim. corollario]
    Per ogni \(\mathcal{F}\) prefascio, sia \(\eta_\mathcal{F} : \mathcal{F} \to
    \mathcal{F}^{\mathcal{A}}\) localmente biunivoca con
    \(\mathcal{F}^{\mathcal{A}}\) fascio. Allora per ogni \(\varphi  :
    \mathcal{F} \to \mathcal{G}\) esiste unico \(\varphi ^{\mathcal{A}} :
    \mathcal{F} ^{\mathcal{A}} \to \mathcal{G}^{\mathcal{A}}\) tale che il
    seguente diagramma commuta (per la proposizione). Allora si conclude
    facilmente (esercizio).
\[\begin{tikzcd}
	{\mathcal{F}} & {\mathcal{F}^\mathcal{A}} \\
	{\mathcal{G}} & {\mathcal{G}^\mathcal{A}}
	\arrow["{\eta_\mathcal{F}}", from=1-1, to=1-2]
	\arrow["\varphi"', from=1-1, to=2-1]
	\arrow["{\varphi^\mathcal{A}}"', from=1-2, to=2-2]
	\arrow["{\eta_\mathcal{G}}", from=2-1, to=2-2]
\end{tikzcd}\]
\end{proof}

\begin{proof}[Idea dim. della proposizione]
    Per ogni \(\mathcal{F}\) prefascio, sia \(\mathcal{F}^{s}\) (``prefascio
    separato associato a \(\mathcal{F}\)'') definito come segue:
    \[
      \forall U \subseteq X \text{ aperto, }\mathcal{F}^{s}{(U)} :=
      \mathcal{F}{(U)} /_\sim 
    \]
    con \(\sigma \sim \sigma' \in \mathcal{F}{(U)} \iff \exists \{U_{i}\}_{i \in
    I} \) ricoprimento aperto di \(U\) tale che \(\sigma|_{U_{i}}  =
    \sigma'|_{U_{i}} \). Allora \(\sim \) è una relazione di equivalenza (per la
    transitività si usa che dati 2 ricoprimenti aperti \(\{U_{i}\}_{i \in I} \)
    e \(\{V_{j}\}_{j \in J} \) di \(U\) esiste un ``raffinamento'' comune
    \(\{W_k\}_{k \in K} \) ricoprimento aperto di \(U\), cioè tale che \(\forall
    k \in K\) esista \(i \in I\) e \(j \in J\) tali che \(W_k \subseteq U_{i},
    V_{j} \) e questo è vero poiché si può prendere il ricoprimento dato da
    \(\{U_{i} \cap V_{j}\}_{i \in I, j \in J}\).

    Allora per ogni \(\varphi : \mathcal{F} \to \mathcal{G}\), abbiamo che
    \(\varphi_U^{s} : \mathcal{F}^{s}_{{(U)}} \to \mathcal{G}^{s}_{{(U)}} \)
    data da \([\sigma] \mapsto [\varphi_U{(\sigma)}]\). In tal modo si ottiene
    un prefascio separato \(\mathcal{F}^{s}\) (esercizio). Il morfismo naturale
    \(\mathcal{F} \to \mathcal{F}^{s}\) (di componenti le proiezioni al
    quoziente \(\mathcal{F}{(U)} \to \mathcal{F}^{s}{(U)}\) è suriettivo e
    localmente iniettivo. 

    Allora definisco \(\mathcal{F}^{\mathcal{A}}\) a partire da
    \(\mathcal{F}^{s}\) nel modo seguente:
    \[
      \mathcal{F}^{\mathcal{A}}{(U)} := \{{(\{U_{i}\}_{i \in I} ,
      \{\sigma_{i}\}_{i \in I}  )} : \{U_{i}\}_{i \in I} \text{ ric.
aperto di \(U\), } \sigma_{i} \in \mathcal{F}^{s}{(U_{i})}\,\,\forall i \in I \}/_\sim 
    \] % TODO agigungere incollamento
    dove \({(\{U_{i}\} _{i \in I} , \{\sigma_{i}\} _{i \in I} )} \sim
    {(\{V_{j}\} _{j \in J} , \{\tau_{j}\} _{j \in J} )}\) se e solo se
    \(\sigma_{i}|_{U_{i} \cap V_{j}} = \tau_{j}|_{U_{i} \cap V_{j}} \), per ogni
    \(i \in I\) e \(j \in J\). Allora \(\sim \) è relazione d'equivalenza (dove
    per la transitività si usa che \(\mathcal{F}^{s}\) è separato)

    Allora \(\mathcal{F}^{\mathcal{A}}\) è prefascio poiché \(\forall U
    \subseteq V \subseteq X  \) aperti, \[[{(\{V_{i}\}_{i \in I}  ,
        \{\sigma_{i}\}_{i \in I} )}] \in
    \mathcal{F}^{\mathcal{A}}{(V)} \mapsto [{(\{U_{i}\}_{i \in I} ,
    \{\sigma_{i}\} _{i \in I} )}] \in
\mathcal{F}^{\mathcal{A}}{(U)}\] e si verifica che \(\mathcal{F}^{\mathcal{A}}\)
è un fascio.

    La mappa naturale \(\mathcal{F}^{s} \to \mathcal{F}^{\mathcal{A}}\), data
    per ogni \(U \in \mathrm{Open}{(X)}\) da \(\sigma \mapsto [{(\{U\} ,
    \{\sigma\} )}]\) è un morfismo di prefasci localmente suriettivo (per
    costruzione) e anche iniettivo (perché \(\mathcal{F}^{s}\) è separto).
    Ottengo allora \(\eta\) come composizione di \(\mathcal{F} \to
    \mathcal{F}^{s} \to \mathcal{F}^{\mathcal{A}}\). Allora \(\eta\) è
    localmente biunivoca perché lo sono entrambe.
\end{proof}

    \begin{lemma}{}
        Sia \(0 \to \mathcal{F}' \overset{\varphi }{\to} \mathcal{F}
        \overset{\psi}{\to} \mathcal{F}'' \to 0\) esatta in
        \(O_X\mathtt{-PMod}\), cioè \(\forall U \in \mathrm{Open}{(X)}\), \(0
        \to \mathcal{F}'{(U)} \overset{\varphi _U}{\to} \mathcal{F}{(U)}
        \overset{\psi_U}{\to} \mathcal{F}''{(U)} \to 0\) è esatta in
        \(O_X{(U)}\mathtt{-Mod}\) o in \(\mathtt{Ab}\).

        Allora:
\begin{enumerate}[label = \arabic*)]
    \item Se \(\mathcal{F}\) è separato, anche \(\mathcal{F}'\) è separato
    \item Se \(\mathcal{F}\) è separato e \(\mathcal{F}'\) è un fascio, allora
        \(\mathcal{F}''\) è separato
    \item Se \(\mathcal{F}\) è fascio e \(\mathcal{F}''\) è seaparato, allora
        \(\mathcal{F}'\) è fascio
\end{enumerate}
    \end{lemma}

\begin{proof}{}
\begin{enumerate}[label = \arabic*)]
    \item Sia \(U \in \mathrm{Open}{(X)}\) e \(\sigma' \in \mathcal{F}'{(U)}\)
        tali che esista \(\{U_{i}\}_{i \in I}  \) ricoprimento aperto di \(U\)
        con \(\sigma'|_{U_{i}} = 0\). Devo dimostrare che \(\sigma' = 0\).
        Sappiamo che 
        \[
            0 = \varphi {(\sigma')}|_{U_{i}}  = \varphi {(\sigma')}|_{U_{i}}
          \implies \varphi {(\sigma')} = 0 \text{ (perché \(\mathcal{F}\) è
          separato)}
        \]
        poiché \(\varphi \) è iniettivo, allora \(\sigma' = 0\).
    \item Sia \(U \in \mathrm{Open}{(X)}\), \(\sigma'' \in \mathcal{F}''{(U)}\)
        tale che esista un ricoprimento aperto \(\{U_{i}\} _{i \in I} \) di
        \(U\) con \(\sigma''|_{U_{i}} = 0\) per ogni \(i \in I\). Devo
        dimostrare che \(\sigma'' = 0\). Poiché \(\varphi \) è suriettivo,
        \(\exists \sigma \in \mathcal{F}{(U)}\) tale che \(\sigma'' = \psi
        {(\sigma)}\). Allora \(\psi{(\sigma|_{U_{i}})} = \psi{(\sigma)}|_{U_{i}}
        = \sigma''|_{U_{i}} = 0\). Per l'esattezza, \(\exists !\, \sigma'_i \in
        \mathcal{F}'{(U_{i})}\) tale che \(\sigma|_{U_{i}}  = \varphi
        {(\sigma'_i)}\). 
        Ora abbiamo che per ogni \(i, j \in I\), abbiamo che
        \[
          \varphi {(\sigma'_i|_{U_{i, j} })} = \varphi {(\sigma'_i)}|_{U_{i,j}}
          = {(\sigma|_{U_{i}} )}|_{U_{i,j}} = \sigma|_{U_{i,j}} = \varphi
          {(\sigma'_j|_{U_{i,j}} )}
        \]
        allora, poiché \(\varphi \) è iniettivo, \(\sigma'_j|_{U_{i,j}} =
        \sigma'_j|_{U_{i,j}}\).

        Infine poiché \(\mathcal{F}\) è fascio, esiste unico \(\sigma' \in
        \mathcal{F}'{(U)}\) tale che \(\sigma'_i = \sigma'|_{U_{i}} \) per ogni
        \(i \in I\). Allora
        \[
          \varphi {(\sigma)}|_{U_{i}} =\varphi {(\sigma'|_{U_{i}} )} = \varphi {(\sigma'_i)} =
          \sigma|_{U_{i}} \implies \varphi {(\sigma')} = \sigma
        \]
        perché \(\mathcal{F}\) è separato, da cui \(\sigma'' = \psi{(\sigma)} =
        \psi{(\varphi {(\sigma')})} = 0\) 
    \item So già che \(\mathcal{F}'\) è separato per il punto 1). Resta da
        dimostrare che dato \(\{U_{i}\} _{i \in I} \) ricoprimento aperto di
        \(U\) e \(\sigma'_i \in \mathcal{F}'{(U_{i})}\) tali che
        \(\sigma'_i|_{U_{i,j}} \) per ogni \(i , j \in I\), allora esiste
        \(\sigma' \in \mathcal{F}'{(U)}\) tale che \(\sigma'_i =
        \sigma'|_{U_{i}} \) per ogni \(i \in I\). Ma allora
        \[
          \varphi {(\sigma'_i)}|_{U_{i,j}} = \varphi {(\sigma'_i|_{U_{i,j}} )} =
          \varphi {(\sigma'_j|_{U_{i,j}} )} = \varphi {(\sigma'_j)}|_{U_{i,j}} 
        \]
        Poiché \(\mathcal{F}\) è un fascio, allora esiste unico \(\sigma \in
        \mathcal{F}{(U)}\) tale che \(\sigma|_{U_{i}} = \varphi {(\sigma'_i)}\)
        per ogni \( i \in I\). Allora
        \[
          \psi {(\sigma)}|_{U_{i}} = \psi {(\sigma|_{U_{i}} )} =
          \psi{(\varphi {(\sigma'_i)})} = 0 \implies \psi{(\sigma)} = 0
        \]
        poiché \(\mathcal{F}''\) è separato, allora per l'esattezza esiste unico
        \(\sigma' \in \mathcal{F}'{(U)}\) tale che \(\varphi {(\sigma')} =
        \sigma\). Segue che \(\varphi {(\sigma'|_{U_{i}} )} = \varphi
        {(\sigma')}|_{U_{i}} = \sigma|_{U_{i}}  = \varphi {(\sigma'_i)}\). Segue
        che \(\sigma'_i = \sigma'|_{U_{i}} \) per ogni \(i \in I\) perché
        \(\varphi \) è iniettivo.
\end{enumerate}
\end{proof}

\begin{theorem}{}
    Sia \({(X, \mathcal{O_X})}\) uno spazio anellato. Allora
    \(O_X\mathtt{-Mod}\) è una categoria abeliana.
    \begin{note}[zione]
        \(\mathrm{(Co)ker}_{\mathcal{P}} \) indica un (co)nucleo in
        \(O_X\mathtt{-PMod}\), mentre i corrispettivi in \(O_X\mathtt{-Mod}\)
        saranno indicati con \(\mathrm{(Co)ker}\).
    \end{note}

    Inoltre \(\forall \varphi : \mathcal{F} \to \mathcal{G}\) in
    \(O_X\mathtt{-Mod}\), valgono le seguenti:
\begin{itemize}[label = --]
    \item \(\mathrm{Ker}{(\varphi )} =
    \mathrm{Ker}_\mathcal{P}{(\varphi )}\) (e idem per \(\mathrm{ker}\)).
\item \(\mathrm{Coker}{(\varphi )} = \mathrm{Coker}_\mathcal{P}{(\varphi
    )}^{\mathcal{A}}\) (e \(\mathrm{coker}\varphi =
    \eta_{\mathrm{Coker}_\mathcal{P}{(\varphi)}} \circ
    \mathrm{coker}_{\mathcal{P}}{(\varphi)}\) 
    \item \(\varphi \) è mono in \(O_X\mathtt{-Mod}\) se e solo se \(\varphi \)
        è iniettivo (se e solo se \(\varphi \) è loc. iniettivo \(\iff\)
        \(\varphi \) è mono in \(O_X\mathtt{-PMod}\))
    \item \(\varphi \) è un epi in \(O_X\mathtt{-Mod}\) se e solo se \(\varphi
        \) è localmente suriettivo.
\end{itemize}
\end{theorem}
\begin{proof}{}
    \(O_X\mathtt{-Mod}\) è additivo perché è sottocategoria piena chiusa per per
    (co)prodotti finiti di \(O_X\mathtt{-PMod}\) (abeliana). In generale per
    ogni \(\varphi  : \mathcal{F} \to \mathcal{G}\) in \(O_X\mathtt{-PMod}\) si
    fattorizza in \(O_X\mathtt{-PMod}\)
    % TODO diagramma
    e ci sono successioni esatte in \(O_X\mathtt{-PMod}\) 
    \begin{align*}
      0 \to \mathrm{Ker}_\mathcal{P}{(\varphi )}
      \overset{\mathrm{ker}_\mathcal{P}{(\varphi )}}{\to} \mathcal{F}
      \overset{\mathrm{coim}_\mathcal{P}{(\varphi )}}{\to}
      \mathrm{Im}_\mathcal{P}{(\varphi )} \to 0 \\
      0 \to \mathrm{Im}_\mathcal{P}{(\varphi )}
      \overset{\mathrm{im}_\mathcal{P}{(\varphi )}}{\to} \mathcal{G}
      \overset{\mathrm{coker}_\mathcal{P}{(\varphi )}}{\to}
      \mathrm{Coker}_\mathcal{P}{(\varphi )} \to 0
    \end{align*}
    assumendo \(\varphi \in O_X\mathtt{-Mod}\), per 1)
    \(\mathrm{Im}_\mathcal{P}{(\varphi )}\) è separato, e dunque per 3)
    \(\mathrm{Ker}_\mathcal{P}{(\varphi )}\) è un fascio. Segue che
    \(\mathrm{ker}{(\varphi )} = \mathrm{ker}_\mathcal{P}{(\varphi )}\) e
    \(\mathrm{coker}{(\varphi )} = \eta_{\mathrm{Coker}_\mathcal{P}{(\varphi )}}
    \circ \mathrm{coker}_\mathcal{P}{(\varphi )}\) per la proprietà di
    aggiunzione di \({(-)}^{\mathcal{A}}\) 
\[ \begin{tikzcd}
	\mathcal{F} & \mathcal{G} & \mathrm{Coker}_\mathcal{P}{(\varphi )} &
    \mathrm{Coker}_\mathcal{P}{(\varphi )}^{\mathcal{A}} \\
	& \mathcal{H}
	\arrow["\varphi", from=1-1, to=1-2]
	\arrow["0"', from=1-1, to=2-2]
	\arrow["{\mathrm{coker}_\mathcal{P}(\varphi)}", from=1-2, to=1-3]
	\arrow["{\psi\in O_X\mathtt{-Mod}}", from=1-2, to=2-2]
	\arrow["\eta", from=1-3, to=1-4]
	\arrow["{\varphi'}"{pos=0.3}, curve={height=-6pt}, dashed, from=1-3, to=2-2]
	\arrow["{\varphi''}"{description}, curve={height=-12pt}, dashed, from=1-4, to=2-2]
\end{tikzcd}
\]
  Allora \(\varphi \) è mono se e solo se \(\mathrm{Ker}{(\varphi )} = 0\)   in
  \(O_X\mathtt{-Mod}\). Ma \(\mathrm{Ker}{(\varphi )} =
  \mathrm{Ker}_\mathcal{P}{(\varphi )}\) dunque tale condizione è equivalente a
  \(\varphi \) mono in \(O_X\mathtt{-PMod}\) (dunque \(\iff\) \(\varphi \) è
  iniettivo).

  \textbf{SALTATA LA PARTE \(\varphi \) EPI, NON STO CAPENDO CHE STIAMO
  FACENDO}.

  Sia \(\varphi \) mono, allora \(\varphi =
  \mathrm{ker}_\mathcal{P}{(\mathrm{coker}_\mathcal{P}{(\varphi )})}\). Inoltre
  se \(\varphi \) è mono, \(\mathcal{F} \cong \mathrm{Im}_\mathcal{P}{(\varphi
  )}\) che è allora un fascio. Allora \(\mathrm{Coker}_\mathcal{P} {(\varphi
  )}\) è separato per 2) del lemma precedente. Allora segue che
  \(\eta_{\mathrm{Coker}_\mathcal{P}{(\varphi )}} \) è iniettiva. Allora
  \(\varphi = \mathrm{Ker}_\mathcal{P}{(\mathrm{coker}\varphi )} =
  \mathrm{ker}{(\mathrm{coker})}{(\varphi)}\).

  Inoltre se \(\varphi \) è epi in \(O_X\mathtt{-Mod}\), allora \(\varphi \) è
  localmente suriettiva, da cui (ricordando che \(\varphi =
  \mathrm{im}_\mathcal{P}{(\varphi )} \circ \mathrm{coim}_\mathcal{P}{(\varphi
  )}\)) \(\mathrm{im}_\mathcal{P}{(\varphi )}\) è
  localmente suriettiva, ed è iniettiva. Allora posso assumere
  \(\mathrm{im}_\mathcal{P}{(\varphi )} = \eta_{\mathrm{Im}_\mathcal{P}{(\varphi
  )}} \), per cui
    \begin{align*}
        \mathrm{coim}_\mathcal{P}{(\varphi )}
        &= \mathrm{coker}_\mathcal{P}{(\mathrm{ker}_\mathcal{P}{(\varphi )})}
            = \mathrm{coker}_\mathcal{P}{(\mathrm{ker}{(\varphi )})} \\
        &\implies \varphi
            = \mathrm{im}_\mathcal{P}{(\varphi )} \circ
                \mathrm{coker}_\mathcal{P}{(\mathrm{ker}{(\varphi )})} \\
        &= \eta_{\mathrm{Coker}_\mathcal{P}{(\mathrm{ker} \varphi })} \circ
             \mathrm{coker}_\mathcal{P}{(\mathrm{ker}{(\varphi )})} \\
        &= \mathrm{coker}{(\mathrm{ker}{(\varphi )})}
    \end{align*}
\end{proof}

\section{Sottoprefasci e coomologia locale}

\begin{definition}{Sotto(pre)fascio}
    Sia \(\mathcal{F}\) un (pre)fascio. \(\mathcal{F}' \subseteq \mathcal{F} \)
    è un \emph{sottoprefascio} se \(\forall U \in \mathrm{Open}{(X)}\),
    \(\mathcal{F}'{(U)} \subseteq \mathcal{F}{(U)} \) e \(\forall \sigma \in
    \mathcal{F}'{(U)}\), \(\forall V \subseteq U \) aperto, \(\sigma|_{V} \in
    \mathcal{F}^{\prime}{(V)}\)
\end{definition}
    
    Se \(\mathcal{F}\) è separato, allora \(\mathcal{F}^{\prime}\) è separato.

    Se \(\mathcal{F}\) è un fascio, allora \(\mathcal{F}'^{a}\) si descrive come
    \(\mathcal{F}'^{a}{(U)} := \{\sigma \in \mathcal{F}{(U)}: \exists
    \{U_{i}\}_{i \in I}\\ \text{ ric. aperto di \(U\) t.c. }\sigma|_{U_{i}} \in
\mathcal{F}'{(U_{i})}, \, \forall i \in I\}\) 

In particolare se \(\varphi : \mathcal{F} \to \mathcal{G}\) è un morfismo di
\(0_X\mathtt{-Mod}\), allora \(\mathrm{coim}\varphi =
\mathrm{coker}{(ker{(\varphi )})} = \mathrm{coker}{(ker_\mathcal{P}{(\varphi
)})} = \eta \circ \mathrm{coker}_\mathcal{P}{(\mathrm{ker}_\mathcal{P}{(\varphi
)})} = \eta \circ \mathrm{coim}_\mathcal{P}{(\varphi )}\), da cui
\(\mathrm{Im}{(\varphi )} = \mathrm{Im}_\mathcal{P}{(\varphi )}^{a}\). Allora

Allora \(\mathrm{Im}_\mathcal{P}{(\varphi )} \subseteq \mathcal{G} \) è un
sottoprefascio (separato).
\[
  \implies \mathrm{Im}{(\varphi )}{(U)} = \{\sigma \in \mathcal{G }{(U)} :
  \exists \{U_{i}\} _{i \in I} \text{ ric. aperto di \(U\) t.c. }\sigma|_{U_{i}}
\in \mathrm{Im}{(\varphi_{U_{i}}} \, \forall i \in I\} 
\]

Sia \({(X, 0_X)}\) uno spazio anellato. Per ogni \(U \in \mathrm{Open}{(X)}\)
c'è un funtore \(\Gamma{(U, -)} : 0_X\mathtt{-Mod} \to 0_X{(U)}\mathtt{-Mod}\).
Tale funtore è esatto a sinistra perché è composizione dell'inclusione
\(0_X\mathtt{-Mod} \hookrightarrow 0_X\mathtt{-PMod}\) (che è esatto a sinistra
perché aggiunto destro di \({(-)}^{a}\)) con \(\Gamma{(U, -)}:
0_X\mathtt{-PMod} \to 0_X{(U)}\mathtt{-Mod}\), definito da \(\mathcal{F} \mapsto
\mathcal{F}{(U)}\) che è esatto.

\begin{remark}{}
    In generale \(\Gamma{(X, -)}\) non è esatto (a destra), cioè non preserva
    gli epimorfismi. In altre parole esiste \(\varphi  : \mathcal{F} \to
    \mathcal{G}\) in \(0_X\mathtt{-Mod}\) localmente suriettivo ma tale che
    \(\varphi _X : \mathcal{F}{(X)} \to \mathcal{G}{(X)}\) non è suriettivo.
    Il prossimo esempio mostra che questo è vero
\end{remark}

\begin{example}{}
    Suppongo che \(\exists A\) anello tale che \(0_X{(U)} \subseteq A^{U} \) per
    ogni \(U \in \mathrm{Open}{(X)}\). Allora \(\forall Y \subseteq X \) sia
    \(\mathcal{I}_Y \subseteq 0_X \) sottofascio (di \(0_X\)-moduli, detto di
    ``ideali''), definito per ogni \(U \in \mathrm{Open}{(X)}\) da
    \(\mathcal{I}_Y{(U)} := \{\sigma \in 0_X{(U)} : \sigma|_Y = 0\}\). Ottengo
    allora un epimorfismo in \(0_X\mathtt{-Mod}\) \(\pi: 0_X \to 0_X /
    \mathcal{I}_Y\). Cerco un esempio in cui \(\pi_X : 0_X {(X)} \to {(0_X /
    \mathcal{I}_Y)}{(X)}\) non suriettivo. Preso ora \(0_X = A_X \) (fascio
    localmente costante) con \(A \neq 0\), \(X\) connesso e \(Y = \{x_{1},
    x_{2}\} \) con \(x_{1} \neq x_{2}\) punti chiusi.

Allora \(U_{1} = X \sminus \{x_{2}\} \) e \(U_{2} = X \sminus \{x_{1}\}\)
    sono aperti e \(\{U_{1}, U_{2}\} \) è un ricoprimento aperto di \(X\). Sia
    ora \(a_{1}, a_{2} \in A\) con \(a_{1} \neq a_{2}\). Sia \(\sigma_{i} \in
    0_X{(U_{i})}\) definito da \(\sigma_i{(x)} := a_{i}\) per ogni \(x \in
    U_{i}\). Allora \(\sigma_{2}|_{U_{1,2}}, \sigma_{1}|_{U_{1,2}} \in
    \mathcal{I}_Y{(U_{1,2})} = A_X{(U_{1,2} )}\).

    \begin{align*}
      &\implies \overline{\sigma}_i := \pi{(\sigma_i)}\in {(0_X /
      \mathcal{I}_Y)}{(U_{i})} \text{ t.c. } \overline{\sigma}_1 |_{U_{1,2} } =
        \overline{\sigma}_2|_{U_{1,2}} \\&\implies \exists!\, \tau \in {(0_X /
      \mathcal{I}_Y)}{(X)} \text{ t.c. } \tau|_{U_{i}} = \overline{\sigma}_i
    \end{align*}

    Se ora per assurdo \(\exists \sigma \in 0_X{(X)}\) tale che \(\pi{(\sigma)}
    = \tau\), allora \(\sigma{(x)} = a \in A\) per ogni \(x \in X\). Ma
    \[
      \pi{(\sigma|_{U_{i}} )} = \pi{(\sigma)}|_{U_{i}} = \tau|_{U_{i}} =
      \overline{\sigma}_i = \pi{(\sigma_i)} \implies \sigma|_{U_{i}} - \sigma_i
      \in \mathcal{I}_Y{(U_{i})}
    \]
    dunque \(a = \sigma{(x_{i})} = \sigma_i{(x_{i})} = a_{i}\) che è assurdo.
\end{example}

\begin{note}[zione]
    Per ogni \(\mathcal{F}, \mathcal{G} \in 0_X\mathtt{-Mod}\) scrivo
    \(\mathrm{Hom}_X{(\mathcal{F}, \mathcal{G})}\) o \(\mathrm{Hom}_{0_X}
    {(\mathcal{F}, \mathcal{G})}\) invece di \(0_X\mathtt{-Mod}{(\mathcal{F},
    \mathcal{G})}\).

    \(\mathrm{Hom}_X{(0_X, \mathcal{F})}\) è in modo naturale un
    \(0_X{(X)}\)-modulo.
\end{note}

\begin{remark}{}
    Il funtore \(\mathrm{Hom}_X{(0_X, -)} : 0_X\mathtt{-Mod} \to
    0_X{(X)}\mathtt{-Mod}\) è isomorfo a \(\Gamma{(X, -)}\).
    Infatti per ogni \(\mathcal{F} \in 0_X\mathtt{-Mod}\),
    \(\mathrm{Hom}_X{(0_X, \mathcal{F})} \to \mathcal{F}{(X)}\) definito da
    \(\varphi  \mapsto \varphi _X{(1)}\) ha inversa \(\mathcal{F}{(X)} \to
    \mathrm{Hom}_X{(0_X, \mathcal{F})}\) dato da \(\sigma \mapsto
    \widehat{\sigma}\), con \(\widehat{\sigma}_U{(\alpha)} := \alpha \sigma|_U
    \in \mathcal{F}{(U)}\) (\(\alpha \in 0_X{(U)}\)).

    Quindi in generale \(\mathrm{Hom}_X{(0_X,-)}\) non è esatto, cioè \(0_X\)
    non è oggetto proiettivo di \(0_X\mathtt{-Mod}\).
    È possibile dimostrare che non ci sono abbastanza proiettivi in
    \(0_X\mathtt{-Mod}\). D'altra parte ci sono sempre abbastanza iniettivi
    (vedremo).
\end{remark}

\begin{definition}{Coomologia di fasci}
    Per ogni \(n \in \mathbb{N}\), \(H^{n}{(X, -)} := R^{n}\Gamma{(X, -)} :
    0_X\mathtt{-Mod} \to 0_X{(X)}\mathtt{-Mod}\).

    Quindi per ogni successione esatta \(0 \to \mathcal{F}' \to \mathcal{F} \to
    \mathcal{F}'' \to 0\) in \(0_X\mathtt{-Mod}\) ottengo una successione esatta
    \[
      0 \to H^{0}{(X, \mathcal{F}')} \to H^{0}{(X, \mathcal{F})} \to H^{0}{(X,
      \mathcal{F}'')} \to H^{1}{(X, \mathcal{F}')} \to \dots
    \]
\end{definition}

\begin{proposition}{}
    Sia \({(X, O_X)}\) uno spazio anellato. Allora in \(O_X\mathtt{-Mod}\) ci
    sono prodotti e coprodotti (i.e. è completa e cocompleta).
\end{proposition}
\begin{proof}{}
    È chiaro per \(O_X\mathtt{-PMod}\), infatti
    \[
      {(\prod_{\lambda \in \Lambda} \mathcal{F}_\lambda)} {(U)} = \prod_{\lambda
      \in \Lambda} \mathcal{F}_\lambda{(U)} \in O_X{(U)}\mathtt{-Mod} 
    \]
    e analogamente per \(\coprod_{\lambda \in \Lambda} \mathcal{F}_\lambda\) 

    Se \(\mathcal{F}_\lambda\) sono fasci, allora \(prod_{\lambda \in \Lambda}
    \mathcal{F}_\lambda\) è un fascio (esercizio). %TODO eser
    Per il coprodotto, risulta che \(\coprod_{\lambda \in \Lambda}
    \mathcal{F}_\lambda = {\left( \coprod^{\mathcal{P}}_{\lambda \in \Lambda}
    \mathcal{F}_\lambda \right)} ^{a}\), poiché \({(-)}^{a}\) preserva i
    colimiti e \(\mathcal{F}_\lambda^{a}\cong \mathcal{F}_\lambda\) se
    \(\mathcal{F}_\lambda\) è un fascio.
\end{proof}

\begin{definition}{Spiga}
    Sia \(\mathcal{F}:\mathrm{Open}{(X)}^{op} \to \mathcal{C} \in \mathtt{PSh}{(X, \mathcal{C})}\). Allora la
    \emph{spiga} di \(\mathcal{F}\) in \(x \in X\) è 
    \[
      \mathcal{F}_x := \mathrm{colim} \,\mathcal{F}|_{\mathrm{Open}_x{(X)}^{op}} \in
      \mathcal{C}
    \]
    dove \(\mathrm{Open}_x{(X)} := \{U \in \mathrm{Open}{(X)} : x \in U\} \).
    Allora \(\mathrm{Open}_x{(X)}^{op}\) è un \emph{insieme filtrante}, cioè
    della forma \({(I, \le )}\) insieme (pre)ordinato tale che \(\forall i, j
    \in I\), \(\exists k \in I\) tale che \(i , j \le k\) (infatti dati \(U, V
    \in \mathrm{Open}_x{(X)}\) esiste \(W \in \mathrm{Open}_x{(X)}\) tale che
    \(W \subseteq U, V \)).
\end{definition}

Sia \(F : {(I, \le )}\to \mathcal{C}\) un funtore con \({(I, \le )}\) filtrante
e \(\mathcal{C}\) una categoria concreta. Allora
\[
  \mathrm{colim}\, F \overset{\text{insieme}}{=} {\left( \coprod_{i \in I} F{(i)}
  \right)} /_\sim 
\]
con \(a \sim b \iff \exists k \in I\) tale che \(i, j \le k\) e \(F{(i\to
k)}{(a)} = F{(j \to k)}{(b)} \in F{(k)}\) con morfismi naturali \(F{(i)} \to
\mathrm{colim}\, F\).

Dunque esplicitamente possiamo avere
\[
  \mathcal{F}_x = {\left( \coprod_{U \in \mathrm{Open}{(X)}} \mathcal{F}{(U)}
  \right)}/_\sim 
\]
con \(\sigma \in \mathcal{F}{(U)} \sim \tau \in \mathcal{F}{(V)}\iff \exists \,W
\in \mathrm{Open}_x{(U \cap V)}\) tale
che \(\sigma|_W = \tau|_W\) e gli elementi di \(\mathcal{F}_x\) vengono detti
\emph{germi} di sezioni di \(\mathcal{F}\) in \(x\).

\begin{note}[zione]
    Data \(\sigma \in \mathcal{F}{(U)}\), con \(U \in \mathrm{Open}_x {(X)}\),
    allora indico con \(\sigma_x = [\sigma] \in  \mathcal{F}_x\).
\end{note}

Se \(\mathcal{F} \in O_X\mathtt{PMod}\), allora \(\mathcal{F}_x\) è
\(O_{X,x}\)-modulo.

\(\varphi : \mathcal{F} \to \mathcal{G}\) morfismo di prefasci induce morfismi
in \(\mathcal{C}\) di \(O_X,x\)-moduli \(\varphi_x : \mathcal{F}_x \to
\mathcal{G}_x\) per ogni \(x \in X\) tale che \(\forall U \in
\mathrm{Open}_x{(X)}\) e \(\forall \sigma \in \mathcal{F}{(U)}\),
\(\mathcal{F}_x{(\sigma_x)} = \varphi _U{(\sigma)}_x\) 

\begin{lemma}{}\label{lem:inj}
    Sia \(\varphi : \mathcal{F} \to \mathcal{G}\) morfismo di prefasci. Allora
    \(\varphi \) è localmente \emph{iniettivo}/\emph{suriettivo} se e solo se
    \(\varphi_x : \mathcal{F}_x \to \mathcal{G}_x\) è
    \emph{iniettivo}/\emph{suriettivo} per ogni \(x \in X\).

    In particolare \(\eta_\mathcal{F} : \mathcal{F} \to \mathcal{F}^{a}\) induce
    un isomorfismmo sulle spighe.
\end{lemma}

\begin{remark}{}
    Una successione in \(O_X\mathtt{-Mod}\) è esatta se e solo se lo è sulle
    spighe.
\end{remark}

\begin{proof}[Dimostrazione lemma (\emph{iniettivo})] \( \) 
    \begin{itemize}
        \item[\(\implies \)] Sia \(\sigma_x, \tau_x \in \mathcal{F}_x\), con
            \(\sigma \in \mathcal{F}{(U)}\), \(\tau \in \mathcal{F}{(V)}\) e
            \(U, V \in \mathrm{Open}_x{(X)}\) tali che \(\varphi_x{(\sigma_x)} =
            \varphi _x{(\tau_x)}\). Allora devo dimostrare \(\sigma_x =
            \tau_x\).

            Sappiamo già che \(\varphi {(\sigma)}_x = \varphi {(\tau)}_x\).
            Allora esiste \(W \in \mathrm{Open}_x{(U \cap V)}\) tale che
            \(\varphi {(\sigma)}|_W = \varphi {(\tau)}|_W\) ma allora \(\varphi
            {(\sigma |_W)} = \varphi {(\tau|_W)}\). Poiché \(\varphi \) è
            localmente iniettivo, allora esiste un \(W' \in
            \mathrm{Open}_x{(W)}\) tale che \({(\sigma|_W)}|_{W^{\prime}} =
        {(\tau|_W)}|_{W^{\prime}} \) da cui \(\sigma |_{W'} = \tau|_{W^{\prime}}
        \), per cui \(\sigma_x = \tau_x\) 
        \item[\(\impliedby \)] Sia \(U \in \mathrm{Open}{(X)}\) e siano
            \(\sigma, \tau \in \mathcal{F}{(U)}\) tali che \(\varphi {(\sigma)} =
            \varphi {(\tau)}\). Per ogni \(x \in U\), abbiamo che \(\varphi
            {(\sigma)}_x = \varphi {(\tau)}_x\)  da cui \(\varphi _x{(\sigma_x)}
            = \varphi _x{(\tau_x)}\) per cui \(\sigma_x = \tau_x\) poiché
            \(\varphi _x\) è iniettiva. Allora \(\exists \, U_x \in
            \mathrm{Open}_x{(U)}\) tale che \(\sigma|_{U_x}  = \tau_{U_x}\).

            Adesso però \(\{U_x\}_{x \in U}  \) è un ricoprimento aperto di
            \(U\) tale che \(\sigma|_{U_x} = \tau|_{U_x} \) per ogni \(x \in
            U\), e quindi \(\varphi \) è localmente iniettiva.
    \end{itemize}
\end{proof}

\begin{proposition}{}
    Sia \({(X, O_X)}\) uno spazio anellato. Allora \(O_X\mathtt{-Mod}\) ha
    abbastanza iniettivi.
\end{proposition}
\begin{proof}{}
    Per ogni \(x \in X\), 
    \begin{align*}
        O_X\mathtt{-Mod} &\longrightarrow O_{X,x}\mathtt{-Mod} \\
         \mathcal{F}&\longmapsto \mathcal{F}_x
    \end{align*}
    è aggiunto sinistro di 
    \begin{align*}
        O_{X,x} \mathtt{-Mod}  &\longrightarrow O_X\mathtt{-Mod} \\
        M &\longmapsto \widehat{M}
    \end{align*}
    con \(\widehat{M}{(U)} = M {(x \in U)}\) con restrizione \(1_M\) quando
    possibile.

    Questo sono aggiunti perché \(\forall \mathcal{F} \in O_X\mathtt{-Mod}\) e
    \(\forall M \in O_{X,x}\mathtt{-Mod}\), 
    \[
      \mathrm{Hom}_{O_{X,x} } {(\mathcal{F}_x, M)} \cong
      \mathrm{Hom}_X{(\mathcal{F}, \widehat{M})} \ni \varphi 
    \]
    dato da morfismi di \(O_X{(U)}\)-moduli \(\varphi {(U)} \to \widehat{M}{(U)}
    = M\) per ogni \(U \in \mathrm{Open}_x{(X)}\) tale che \(\forall U \subseteq
    V \in \mathrm{Open}_x{(X)}\). Per la proprietà universale del colimite,
    \(\varphi _U{(\sigma|_V)} = \varphi_V{(\sigma)}\) per ogni \(\sigma \in
    \mathcal{F}{(V)}\).

    L'unità \(\eta\) di questa aggiunzione è tale che \(\forall \mathcal{F} \in
    O_X\mathtt{-Mod}\) e \(\forall x \in X\)
    \[
      \eta_\mathcal{F}{(x)} = \eta_x : \mathcal{F} \to \widehat{\mathcal{F}_x} 
    \]
    induce un isomorfismo sulla spiega in \(x\), \(\mathcal{F}_x \to {\left(
    \widehat{\mathcal{F}_x} \right)}_x\) e ottengo un morfismo in
    \(O_X\mathtt{-Mod}\) \(\mathcal{F} \to \prod_{x \in X}
    \widehat{\mathcal{F}_x}\) indotto dalle \(\eta{(x)}\) per ogni \(x \in X\)
    che è iniettiva su ogni spiga, quindi è (localmente) iniettiva per il
    lemma~\ref{lem:inj}.

    Se dimostro che esiste un monomorfismo \(\varphi {(x)} : \mathcal{F}_x \to
    \mathcal{I}_x\) con \(\mathcal{I}_x\) iniettivo, ottengo
    \[
      \prod_{x \in X} \varphi {(x)} : \prod_{x \in X} \widehat{\mathcal{F}_x}
      \to \prod_{x \in X} \mathcal{I}_x
    \]
    iniettivo con \(\prod_{x \in X} \mathcal{I}_x\) iniettivo.

Più in generale trovo monomorfismo \(\widehat{\mathcal{M}} \to \mathcal{I}\)
    con \(\mathcal{I}\) iniettivo per ogni \(M \in O_{X, x}\mathtt{-Mod}\).

    Esiste \(f : M \to I\) mono in \(O_{X, x} \mathtt{-Mod}\) con \(I\)
    iniettivo. Allora \(\widehat{f} : \widehat{M} \to \widehat{I}\) in
    \(O_x\mathtt{-Mod}\). Poiché \(\widehat{{(-)}}\) è aggiunto destro, allora
    preserva i mono e dunque \(\widehat{f}\) è mono.

    Resta da dimostrare che \(\widehat{I}\) è iniettivo in \(O_X\mathtt{-Mod}\)
    cioè \(\mathrm{Hom}_X{(-, \widehat{I})} \cong\mathrm{Hom}_{O_{X,x} {(-_x,
    I)}} \) manda mono di \(O_X\mathtt{-Mod}\) in epi (di \(\mathtt{Ab}\)).
    Questo è vero perché \(-_x\) preserva i mono per il lemma~\ref{lem:inj} e
    \(\mathrm{Hom}_{O_{X,x} } {(-, I)}\) è esatto perché \(I\) è iniettivo.
    
\end{proof}

\begin{example}{}
    Si può dimostrare che \(\forall A\) anello, \(H^{n}{(X, -)} :
    A_x\mathtt{-Mod} \to A\mathtt{-Mod}\) conincidono con la coomoologia
    singolare se \(X\) è localmente contraibile.
\end{example}

Volendo definire \(\mathrm{Hom}\) in modo che sia un (pre)fascio, possiamo
farlo. Siano \(\mathcal{F}, \mathcal{G} \in O_X\mathtt{-PMod}\). Allora
\(\mathrm{Hom}_X{(\mathcal{F}, \mathcal{G})} \in \mathbb{Z}_X\mathtt{-Mod}\)
(anche \(O_X\mathtt{-Mod}\) se \(O_X\) è commutativo) definito \(\forall U \in
\mathrm{Open}{(X)}\) da 
\[
  \mathrm{Hom}_X{(\mathcal{F}, \mathcal{G})}{(U)} :=
  \mathrm{Hom}_U{(\mathcal{F}|_U, \mathcal{G}|_U)} \in \mathtt{Ab}
\]
dove \(\mathcal{F}|_U\) è prefascio su \(U\) definito per ogni \(V \in
\mathcal{Open}{(U)}\) da \(\mathcal{F}|_U{(V)} = \mathcal{F}{(V)}\). 
Se \(\mathcal{F}, \mathcal{G}\) sono fasci, allora
\(\mathrm{Hom}_X{(\mathcal{F}, \mathcal{G})}\) è un fascio, e si ottiene
dunque un funtore \(\mathcal{Hom}_X{(-, =)} : O_X\mathtt{-Mod}^{op} \times
O_X\mathtt{-Mod} \to \mathbb{Z}_X\mathtt{-Mod}\) esatto a sinistra in ciascun
argomento e si possono considerare i funtori derivati \(\mathtt{Ext}_X^{i}{(-,
=)}\).

Inoltre dati \(\mathcal{F} \in \mathtt{PMod-}O_X\) e \(\mathcal{G} \in
O_X\mathtt{-PMod}\), si definisce anche 
\[
{(\mathcal{F} \otimes_X^{\mathcal{P}} \mathcal{G})} {(U)} = \mathcal{F}{(U)}
  \otimes_{O_X{(U)}} \mathcal{G}{(U)}
\]
per cui \(\mathcal{F} \otimes_X^{\mathcal{P}} \mathcal{G} \in
\mathbb{Z}_X\mathtt{-PMod}\). Se \(\mathcal{F}, \mathcal{G}\) sono fasci, allora
si può prendere \(\mathcal{F} \otimes_X \mathcal{G} := {\left( \mathcal{F}
\otimes_X^{\mathcal{P}} \mathcal{G}\right)}^{a}\) e si ottiene un funtore 
\[
  - \otimes_X = : \mathtt{Mod-}O_X \times O_X\mathtt{-Mod} \to
  \mathbb{Z}_X\mathtt{-Mod}
\]
esatto a destra in ciascun argomento, che si può derivare usando gli
\(O_X\)-moduli piatti ottenendo \(\mathrm{Tor}_i^{X}{(-, =)}\).

\begin{remark}{}
    Sia \(f : X \to Y\) continua. Allora essa induce 
    \begin{align*}
        f_*: \mathtt{(P)Sh}{(X, \mathcal{C})} &\longrightarrow
        \mathtt{(P)Sh}{(Y, \mathcal{C})} \\
        \mathcal{F} &\longmapsto f_*(\mathcal{F}){(V)} := \mathcal{F}{(f^{-1}{(V)})}
    \end{align*}

    e varianti per i fasci di moduli. Ci sono anche funtori \(\mathtt{(P)Sh}{(Y,
    \mathcal{C})} \to \mathtt{(P)Sh}{(X, \mathcal{C})}\).
\end{remark}

\begin{definition}{}
    \(\mathcal{F} \in O_X\mathtt{-Mod}\) è (localmente) di
    \emph{tipo}/\emph{presentazione}
    finita se \(\forall x \in X\) esiste \(U \in \mathrm{Open}_x{(X)}\) e una
    successione esatta della forma 
    \[
      O_U^{n} \to \mathcal{F}|_U \to 0 \quad / \quad O_U^{m} \to O_U^{n} \to
      \mathcal{F}|_U \to 0
    \]
    \(\mathcal{F}\) è \emph{coerente} se è di tipo finito e per ogni \(U \in
    \mathrm{Open}{(X)}\) e per ogni \(\mathcal{F}^{\prime} \subseteq \mathcal{F}
    |_U\) sottofascio di tipo finito è di presentazione finita.

    Si ottiene sempre una categoria abeliana \(\mathrm{Coh}{(X)} \subseteq
    O_X\mathtt{-Mod}\). Essere coerente è equivalente a essere di presentazione
    finita se e solo se \(O_X\) è coerente (vero ad esempio per varietà
    complesse)
\end{definition}



